We have presented the Oracle Loop, a complete architecture for runtime
self-regulation in language models using KV-cache geometry as detection
signal, steering input, and training reward---and demonstrated that two
of its three components work, with methodological refinements that
sharpened both.

For detection, independent statistical audit and feature ablation
revealed that the generation-window LOO classifier (0.707/0.758) was
primarily driven by a length-correlated feature
(\texttt{mp\_norm\_per\_token}, $R^2 = 0.129$). This prompted
analysis of the full KV cache (encoding + generation), which identified
\texttt{mp\_spectral\_gap} as a cleaner detection signal (AUROC 0.903,
95\% CI [0.806, 0.977]; FWL-corrected 0.877; $R^2 = 0.068$). The
full-cache spectral gap outperforms a text-only baseline (0.755) in
96.7\% of bootstrap samples, though the confidence interval on the
added value includes zero. These confabulation results extend
established deception detection at AUROC 1.000 across seven model
scales, validated through 50+ experiments and an independent audit of
135 claims.

For steering, the preliminary 5-arm experiment ($n = 7$
confabulations) established that cache injection modifies behavior---the
null control producing character-level identical output on all 100
trials is an unusually strong mechanism verification. A subsequent
formulary study ($n = 68$ confabulations, 12 vectors on
Qwen2.5-7B-Instruct) provides the first statistically significant
correction results: hostile injection corrects 96\% of confabulations
(McNemar $p < 0.0001$) at 2.4\% adverse rate. Therapeutic rankings
between models are uncorrelated ($\rho = 0.119$, $p = 0.71$),
confirming that per-model calibration is required.

Three rounds of adversarial red-teaming---culminating in an independent
statistical audit that identified the norm-feature dependency and
prompted the full-cache refinement---demonstrate that adversarial
methodology is not just desirable but essential for claims in this space.
The detection methodology improved because it was attacked honestly.

The Oracle Loop creates dual-use risks: cache injection could suppress
safety-relevant hedging, and geometric features could enable
surveillance. A companion paper demonstrates that KV-Cloak cache
obfuscation effectively blinds external geometric analysis (spectral gap
degraded from 0.903 to 0.614 under the real KV-Cloak mechanism),
providing defense against adversarial use of these methods while a
defender inside a trusted execution environment retains full detection
capability. We advocate that any deployment include obfuscation as a
mandatory companion.

The gap that remains is training. The Oracle Loop's detection signal is
validated---more robustly after refinement than before. Its steering
mechanism is validated with statistical power. Its attack surface is
characterized and defended. The next step is closing the loop: GRPO
training with geometric reward to produce models that confabulate less
because they have learned to recognize confabulation geometry from the
inside.

\paragraph{Code and Data.}
Detection experiment code, ablation analysis, formulary results, raw
trial data, Cache Integrity Monitor implementation, KV-Cloak obfuscation
analysis, and all red-team logs are publicly available at:
\url{https://github.com/Liberation-Labs-THCoalition/lyra-s-research-}.
Full experiment infrastructure:
\url{https://github.com/Liberation-Labs-THCoalition/Project-Oracle}.
